\documentclass[aps,prl,twocolumn,showpacs,superscriptaddress]{revtex4-2}
\usepackage{amsmath,amssymb,booktabs,xcolor,hyperref,amsthm}
\newtheorem{theorem}{Theorem}

\begin{document}

\title{Appendix KA2: OHC Phase Memory and Morphic Coupling\\
Derivation of the non-local phase imprint amplitude and the
BCT Repetition Theorem}

\author{Michel Robert Cabri\'e}
\affiliation{Independent Researcher, Victoria, Australia}
\email{ZeroFreeParameters@gmail.com}
\date{20 March 2026}

\begin{abstract}
This appendix derives the OHC phase memory amplitude, the morphic 
coupling constant, and the BCT Repetition Theorem from the OHC 
condensate equations, supporting Letter~79 (Morphic Resonance and 
the OHC Condensate). When a Hopf transition occurs at position 
$\mathbf{x}_0$, a phase perturbation $\delta\varphi = \alpha_0 
\times (\ell_P/r)$ propagates at $v_{\rm ph} = 6.78c$ throughout 
the global OHC condensate. After $n$ repetitions of pattern $P$, 
the accumulated phase $\Phi_n = n\alpha_0(\ell_P/r_0)$ reduces the 
energy barrier by $\Delta E_n = \hbar\omega_{\rm OHC}\cdot\alpha_0^2
\cdot n$. This is the BCT Repetition Theorem.
\end{abstract}

\maketitle

\section{OHC Phase Perturbation}

When Hopf transition $H=n\to0$ occurs at $\mathbf{x}_0$, the 
released topological phase propagates as a spherical wave in the 
global OHC condensate:
\begin{equation}
    \delta\varphi(r,t) = \alpha_0 \cdot \frac{\ell_P}{r} 
    \cdot f\!\left(t - \frac{r}{v_{\rm ph}}\right)
\end{equation}
where $\alpha_0 = 0.0074081$, $\ell_P = 1.616\times10^{-35}$~m, 
$v_{\rm ph} = 6.78c$ (Letter~59), and $f$ is a causal response 
function with $f(t<0)=0$. The amplitude falls as $1/r$ but never 
reaches zero --- the condensate is globally connected.

\section{Propagation Speed}

\begin{center}
\begin{tabular}{lll}
\toprule
Distance & Travel time & Note \\
\midrule
1~m (lab) & $4.9\times10^{-10}$~s & Instantaneous \\
$6.4\times10^6$~m (Earth) & $3.1\times10^{-3}$~s & milliseconds \\
$4.4\times10^{26}$~m (universe) & $\sim$45~yr & Global \\
\bottomrule
\end{tabular}
\end{center}

\section{The BCT Repetition Theorem}

\begin{theorem}[BCT Repetition Theorem]
Let topological pattern $P$ be established $n$ times in the OHC. 
The accumulated phase imprint is:
\begin{equation}
    \Phi_n = n \cdot \delta\varphi_1 = n \cdot \alpha_0 \cdot 
    \frac{\ell_P}{r_0}
\end{equation}
The energy barrier for the $(n+1)$-th establishment of $P$ is 
reduced by:
\begin{equation}
    \Delta E_n = \hbar\omega_{\rm OHC} \cdot \alpha_0^2 \cdot n
\end{equation}
Repeated patterns are progressively easier to establish, with 
facilitation linear in $n$. \qed
\end{theorem}

With $\alpha_0^2 \approx 5.5\times10^{-5}$, the effect is weak but 
accumulates linearly with $n$ and is non-zero for all $n\geq1$.

\section{Non-Electromagnetic Character}

The OHC phase propagates through the vacuum condensate. The coupling 
$\alpha_0 = r_{\rm oct}\times r_{\rm tet}/\pi$ involves only void 
geometry --- no electric charge appears. OHC phase memory is therefore 
non-electromagnetic and cannot be shielded by Faraday cages or 
conducting enclosures. This is Prediction~P2 of Letter~79.

\begin{thebibliography}{99}
\bibitem{AppJH} M.~R.~Cabri\'e, \textit{Appendix JH}, Zenodo (2026),
doi:10.5281/zenodo.18975018
\bibitem{L59} M.~R.~Cabri\'e, \textit{BCT Letter 59}, Zenodo (2026).
\bibitem{L79} M.~R.~Cabri\'e, \textit{BCT Letter 79}, Zenodo (2026).
\end{thebibliography}

\end{document}
